% ch02_u04: 原书 p.44--p.52
\SourcePage{44}

说明一下：(31) 中我们把 $f_r$ 看作取值于 $(\mathfrak{g}^*)^{\otimes r}$ 的函数。

我们不需要 [BD94] 中同构 (30) 的显式公式。只提一点：按 [BD94] 的命题 5，
$\mathfrak{g}\otimes K$ 在 $\overline{U}'_k$ 上的伴随作用经 (30) 诱导出
$\mathfrak{g}\otimes K$ 在 $\Omega_k$ 上的如下作用：$a\in\mathfrak{g}\otimes K$
把 $(w_0,\ldots,w_k)\in\Omega_k$ 送到 $(0,w''_1,\ldots,w''_k)$，其中
\[
\begin{aligned}
w''_r&=\frac{1}{(r-1)!}\operatorname{Sym} w'_r,\\
w'_r(t_1,\ldots,t_r)&:=(\operatorname{id}\otimes\ldots\otimes\operatorname{id}
  \otimes\operatorname{ad}_{a(t_r)})w_r(t_1,\ldots,t_r)\\
&\quad-w_{r-1}(t_1,\ldots,t_{r-1})\otimes c\cdot da(t_r).
\end{aligned}
\]
这里 $\operatorname{Sym}$ 表示对称化算子（不含因子 $1/r!$），
$\operatorname{ad}_{a(t_r)}:\mathfrak{g}^*\to\mathfrak{g}^*$ 是余伴随表示中
对应于 $a(t_r)$ 的算子，$c:\mathfrak{g}\to\mathfrak{g}^*$ 是 $\mathfrak{g}$
的双线性形式。

\emph{注.} 设 $c=0$，且 $\mathfrak{g}$ 交换。则
$U'_k/I_{n,k}=\bigoplus_{r=0}^k\operatorname{Sym}^r(\mathfrak{g}\otimes K/\mathfrak{m}^n)$，
$\Omega_{n,k}=\bigoplus_{r=0}^k\overline{\operatorname{Sym}}^r
(\mathfrak{g}^*\otimes\mathfrak{m}^{-n}\omega_{\mathcal{O}})$，
其中 $\overline{\operatorname{Sym}}^r$ 表示完备对称幂。同构
$U'_k/I_{n,k}\xrightarrow{\sim}(\Omega_{n,k})^*$ 就是把
$\operatorname{Sym}(\mathfrak{g}\otimes K/\mathfrak{m}^n)$ 与
$\mathfrak{g}^*\otimes\mathfrak{m}^{-n}\omega_{\mathcal{O}}$ 上的多项式函数空间
等同起来；这个等同在 2.4.1 与 2.4.2 中用过。

\bdp{2.9.5} 按 2.9.4，要证
$f:\bar{\mathfrak{z}}\to\mathfrak{z}_{\mathfrak{g}}(\mathcal{O})$ 满射，
只需证明：任一 $(\mathfrak{g}\otimes\mathcal{O})$-不变连续线性泛函
$l:\Omega_{0,k}\to\mathbb{C}$ 都可延拓为 $(\mathfrak{g}\otimes K)$-不变连续线性泛函
$\Omega_k\to\mathbb{C}$。考虑连续线性算子
\[
T:\Omega_k\to\mathbb{C}((\zeta))\widehat{\otimes}\Omega_{0,k}
=\Bigl\{\sum_{n=-\infty}^{\infty}a_n\zeta^n
\Bigm| a_n\in\Omega_{0,k},\ a_n\to 0\ (n\to-\infty)\Bigr\},
\]
它由下式定义：
\begin{equation}\tag{32}
T(w_0,\ldots,w_k)=(\widehat{w}_0,\ldots,\widehat{w}_k),\qquad
\widehat{w}_r=w_r(\zeta+t_1,\ldots,\zeta+t_r),
\end{equation}
其中 $w_r(\zeta+t_1,\ldots,\zeta+t_r)$ 看作
\[
\Delta^{-1}\mathbb{C}((\zeta))[[t_1,\ldots,t_n]]\,dt_1\ldots dt_n
=\mathbb{C}((\zeta))\widehat{\otimes}\Delta^{-1}\mathbb{C}[[t_1,\ldots,t_n]]
\,dt_1\ldots dt_n
\]
的元素。
% [Tutor Note: 原书 p.44 把 t_r 写成 t_n，此处照录。]

\SourcePage{45}

\[
\Delta:=\prod_{1\le i<j\le r}(t_i-t_j)^2.
\]
设 $l\in(\Omega_{0,k})^*$。令 $\bar{l}:\Omega_k\to\mathbb{C}((\zeta))$ 为
$T:\Omega_k\to\mathbb{C}((\zeta))\widehat{\otimes}\Omega_{0,k}$ 与
$\operatorname{id}\otimes l:\mathbb{C}((\zeta))\widehat{\otimes}\Omega_{0,k}
\to\mathbb{C}((\zeta))$ 的合成。把 $\bar{l}$ 写成 $\sum_i l_i\zeta^i$，
$l_i\in(\Omega_k)^*$。若 $l$ 是 $\mathfrak{g}\otimes\mathcal{O}$-不变的，
则各泛函 $l_i$ 是 $\mathfrak{g}\otimes K$-不变的。此外
$l_0|_{\Omega_{0,k}}=l$。

\emph{注.} 设 $G$ 是代数群，且 $\operatorname{Lie}G=\mathfrak{g}$。
则 $G(K)$ 作用在 $\mathfrak{g}\otimes K$ 的中心扩张上（见 (19)），
从而作用在 $\overline{U}'_k$ 上；此外 $G(\mathcal{O})$ 作用在 $U'_k/I_{n,k}$ 上。
因此 $G(K)$ 作用在 $\Omega_k$ 上，$G(\mathcal{O})$ 作用在 $\Omega_{n,k}$ 上。
在上述情形下，若 $l$ 是 $G(\mathcal{O})$-不变的，则各 $l_i$ 是 $G(K)$-不变的
（$G(K)$ 在 $\Omega_k$ 上的作用见 [BD94] 的公式 (24)）。注意：即使 $G$ 连通，
$G(K)$ 也未必连通。因此 $G(K)$-不变性不能由 $(\mathfrak{g}\otimes K)$-不变性
直接推出。

\bdp{2.9.6} 由于 $\bar{l}$ 连续，当 $i\to-\infty$ 时 $l_i\to 0$
（即对每个 $n$，只要 $i$ 足够大就有 $l_{-i}(\Omega_{n,k})=0$）。
于是映射 $l\mapsto\bar{l}$ 可看作从 $U'_k/I_{0,k}$ 到
$W_k:=\bigl\{\sum_{i=-\infty}^{\infty}a_i\zeta^i \bigm| a_i\in\overline{U}'_k,\ 
a_i\to 0\ (i\to-\infty)\bigr\}$ 的映射。这些映射定义一个算子
\begin{equation}\tag{33}
\Phi:\mathit{Vac}'\to W:=\bigcup_k W_k
\end{equation}
其中 $\mathit{Vac}'=U'/I_0$ 是扭真空模。按 2.9.5 的解释，$\Phi$ 诱导映射
\begin{equation}\tag{34}
\mathfrak{z}_{\mathfrak{g}}(\mathcal{O})\to
\mathfrak{z}\widehat{\otimes}\mathbb{C}((\zeta))
:=\Bigl\{\sum_{i=-\infty}^{\infty}a_i\zeta^i
\Bigm| a_i\in\mathfrak{z},\ a_i\to 0\ (i\to-\infty)\Bigr\}.
\end{equation}
可以证明 (34) 是环同态（见 ???）。容易看出：(34) 与投影
$\mathfrak{z}\widehat{\otimes}\mathbb{C}((\zeta))
\to\mathfrak{z}_{\mathfrak{g}}(\mathcal{O})((\zeta))$ 的合成把
$\mathfrak{z}_{\mathfrak{g}}(\mathcal{O})$ 映入
$\mathfrak{z}_{\mathfrak{g}}(\mathcal{O})[[\zeta]]$；并且合成
$\mathfrak{z}_{\mathfrak{g}}(\mathcal{O})\to\mathfrak{z}_{\mathfrak{g}}(\mathcal{O})[[\zeta]]
\xrightarrow{\zeta=0}\mathfrak{z}_{\mathfrak{g}}(\mathcal{O})$ 是恒等映射。
% [Tutor Note: 原书 p.45 此处交叉引用印作「???」，底本如此，保留原样。]

\emph{注.} 设 $G$ 是连通代数群，且 $\operatorname{Lie}G=\mathfrak{g}$。
则 (34) 的像中所有元素都是 $G(K)$-不变的（见 2.9.5 的注）。

\SourcePage{46}

\bdp{2.9.7} 可以证明，(33) 与 Feigin 和 Frenkel 构造的算子
$F:\mathit{Vac}'\to W$ 一致（见 [FF92] 引理 1 的证明）。因此 2.9.5
只是 [FF92] 中一部分内容的一个版本。

[FF92] 中 $F$ 的定义可改写如下。置
$W_k^+:=\overline{U}'_k((\zeta))$；$W_k^-$ 由所有满足「当 $i$ 足够大时
$a_{-i}=0$」的 $\sum_i a_i\zeta^i\in W_k$ 组成。
定义 $W^\pm\subset W$ 为 $W^\pm=\bigcup_k W_k^\pm$。$W^+$ 与 $W^-$
有自然的代数结构，$W$ 有自然的 $(W^+,W^-)$-双模结构
（$W$ 是左 $W^+$-模，也是右 $W^-$-模）。考虑线性映射
$\varphi^\pm:\mathfrak{g}\widehat{\otimes}K\to W^\pm$，使得
\[
\varphi^+(1)=1,\qquad \varphi^-(1)=0,
\]
且对 $a\in\mathfrak{g}((t))=\mathfrak{g}\otimes K
\subset\mathfrak{g}\widehat{\otimes}K$，
\[
\varphi^+(a)=a(t-\zeta)\in\mathfrak{g}((t))((\zeta)),\qquad
\varphi^-(a)=a(t-\zeta)\in\mathfrak{g}((\zeta))((t)).
\]
容易证明 $\varphi^\pm$ 是李代数同态。考虑 $W$ 上的
$\mathfrak{g}\widehat{\otimes}K$-模结构
$a\circ w:=\varphi_+(a)w-w\varphi_-(a)$，
$a\in\mathfrak{g}\widehat{\otimes}K$，$w\in W$。则 $F:\mathit{Vac}'\to W$
是 $\mathfrak{g}\widehat{\otimes}K$-模同态，它把 $\mathit{Vac}'$ 的真空向量
送到 $1\in W$。

\bdp{2.9.8} 说明 (34) 与 2.4.2 中其经典类比的关系。

$\overline{U}'$ 带有标准滤链 $\overline{U}'_k$（见 2.9.4）。它诱导滤链
$\mathfrak{z}_k:=\mathfrak{z}\cap\overline{U}_k$。
% [Tutor Note: 原书 p.46 此处印作 \overline{U}_k（无撇），按 2.9.4 的记号应为 \overline{U}'_k，此处照录。]
我们把 $\operatorname{gr}_k\overline{U}':=
\overline{U}'_k/\overline{U}'_{k-1}$ 与
$\operatorname{Sym}^k(\mathfrak{g}\otimes K)$ 的完备化等同，即与 $k$
次齐次多项式函数 $\mathfrak{g}^*\otimes\omega_K\to\mathbb{C}$ 的空间等同，
其中 $\omega_K:=\omega_{\mathcal{O}}\otimes_{\mathcal{O}}K$（称
$\mathfrak{g}^*\otimes\omega_K$ 上的函数 $f$ 为多项式函数，是指：对每个
$n$，它在 $\mathfrak{g}^*\otimes\mathfrak{m}^{-n}$ 上的限制是多项式，即来自
$\mathfrak{g}^*\otimes(\mathfrak{m}^{-n}/\mathfrak{m}^N)$ 上的多项式函数，
$N$ 依赖于 $n$）。记 $\mathfrak{z}^{cl}$ 为 $\mathfrak{g}^*\otimes\omega_K$ 上
$\mathfrak{g}\otimes K$-不变多项式函数的代数。显然
$\operatorname{gr}\mathfrak{z}$ 在 $\operatorname{gr}\overline{U}'$ 中的像
含于 $\mathfrak{z}^{cl}$。

\SourcePage{47}

$\mathfrak{z}$ 的滤链诱导 $\mathfrak{z}\widehat{\otimes}\mathbb{C}((\zeta))$
上的滤链，且映射 (34) 与滤链相容。我们断言下图交换：
\begin{equation}\tag{35}
\begin{tikzcd}[column sep=4.6em, row sep=2.8em]
\operatorname{gr}_k\mathfrak{z}_{\mathfrak{g}}(\mathcal{O}) \arrow[r] \arrow[d, "\sigma"]
  & \operatorname{gr}_k\mathfrak{z}\widehat{\otimes}\mathbb{C}((\zeta)) \arrow[d] \\
\mathfrak{z}^{cl}_{\mathfrak{g}}(\mathcal{O}) \arrow[r, "\nu"]
  & \mathfrak{z}^{cl}\widehat{\otimes}\mathbb{C}((\zeta))
\end{tikzcd}
\end{equation}
这里上方的箭头由 (34) 诱导，$\mathfrak{z}^{cl}_{\mathfrak{g}}(\mathcal{O})$
的定义见 2.4.1，$\sigma$ 是 1.2.5 中的象征映射，$\nu$ 由下式定义：
\begin{equation}\tag{36}
\begin{gathered}
\nu(f):=h(\zeta),\qquad (h(\zeta))(\varphi):=f(\varphi(\zeta+t))\\
f\in\mathfrak{z}^{cl}_{\mathfrak{g}}(\mathcal{O}),\quad
\varphi\in\mathfrak{g}^*\otimes\omega_K,\quad
\varphi(\zeta+t)\in\mathfrak{g}^*((\zeta))[[t]]dt
=(\mathfrak{g}^*\otimes\omega_{\mathcal{O}})\widehat{\otimes}\mathbb{C}((\zeta)).
\end{gathered}
\end{equation}
这里把 $\mathfrak{z}^{cl}_{\mathfrak{g}}(\mathcal{O})$ 与
$\mathfrak{g}^*\otimes\omega_{\mathcal{O}}$ 上 $\mathfrak{g}\otimes\mathcal{O}$-不变
多项式函数的代数等同（参看 2.4.1）。映射 $\nu$ 在 2.4.2 的注中讨论过。

(35) 的交换性由下图 (37) 的交换性推出：
\begin{equation}\tag{37}
\begin{tikzcd}[column sep=4.6em, row sep=2.8em]
(U'_k/I_{n,k})^{*} \arrow[r, "\sim"] \arrow[d, "\sigma^{*}"]
  & \Omega_{n,k} \arrow[d] \\
(\operatorname{Sym}^k(\mathfrak{g}\otimes K/\mathfrak{g}\otimes\mathfrak{m}^n))^{*}
  \arrow[r, "\sim"]
  & ((\mathfrak{m}^{-n}\omega_{\mathcal{O}})^{\otimes k})^{S_k}
\end{tikzcd}
\end{equation}
这里上方的箭头是 (30) 的对偶，
$\sigma:U'_k/I_{n,k}\to
\operatorname{Sym}^k(\mathfrak{g}\otimes K/\mathfrak{g}\otimes\mathfrak{m}^n)$
是象征映射，右侧竖箭头由 $w\mapsto(0,\ldots,0,w)$ 定义。(37) 的交换性
由 (30) 的定义直接推出；见 [BD94]。

\bdsec{2.10. $T^{*}\operatorname{Bun}_G$ 的几何} 本小节应看作附录；
读者完全可以跳过。

置 $\operatorname{Nilp}=\operatorname{Nilp}(G):=p^{-1}(0)$，其中
$p:T^{*}\operatorname{Bun}_G\to\operatorname{Hitch}(X)$ 是 Hitchin 纤维化
（见 2.2.3）。$\operatorname{Nilp}$ 由 [La87] 与 [La88] 引入，名为整体幂零锥
（若 $\mathcal{F}$ 是 $X$ 上 $G$-丛，且
$\eta\in T^{*}_{\mathcal{F}}\operatorname{Bun}_G
=H^0(X,\mathfrak{g}^{*}_{\mathcal{F}}\otimes\omega_X)$，则
$(\mathcal{F},\eta)\in\operatorname{Nilp}$ 当且仅当 $\eta$ 在
$H^0(X,\mathfrak{g}_{\mathcal{F}}\otimes\omega_X)$ 中的像是幂零的）。

\SourcePage{48}

在 2.10.1 中我们将证明：命题 2.2.4(iii) 容易由等式
\begin{equation}\tag{38}
\dim\operatorname{Nilp}=\dim\operatorname{Bun}_G
\end{equation}
推出。

我们还由 (38) 推出：$\operatorname{Bun}_G$ 在 1.1.1 的意义下是良的。
等式 (38) 由 Faltings 与 Ginzburg 证明；在 $G=PSL_n$ 的特殊情形，
它已由 Laumon 证明。在 2.10.2 中我们对其证明作一些评述。在 2.10.3 中
讨论 $\operatorname{Nilp}$ 的不可约分支集。在 2.10.4 中证明：即使 $X$
的亏格等于 $0$ 或 $1$，$\operatorname{Nilp}$ 也是等维的（若 $g>1$，
这由 2.2.4(iii) 推出）。在 2.10.5 中证明 $\operatorname{Bun}_G$ 在 1.1.1
的意义下是极良的。

我们将用 $\mathfrak{g}$ 上不变的数量积把 $\mathfrak{g}$ 与
$\mathfrak{g}^*$ 等同。

\bdp{2.10.1} 假定 (38) 成立，我们来证明 2.2.4(iii)，并证明
$\operatorname{Bun}_G$ 在 1.1.1 的意义下是良的。设
$U\subset T^{*}\operatorname{Bun}_G$ 是满足
$\dim U\le 2\dim\operatorname{Bun}_G$ 的最大开子叠。(38) 的意思是：
$p:T^{*}\operatorname{Bun}_G\to\operatorname{Hitch}(X)$ 在 $0$ 上的纤维
维数为 $\dim\operatorname{Bun}_G$。由于
$\dim\operatorname{Hitch}(X)=\dim\operatorname{Bun}_G$，这说明
$U\supset p^{-1}(0)$。$U$ 关于 $\mathbb{G}_m$ 在
$T^{*}\operatorname{Bun}_G$ 上的自然作用不变。因此
$U=T^{*}\operatorname{Bun}_G$。于是
$\dim T^{*}\operatorname{Bun}_G\le 2\dim\operatorname{Bun}_G$。
按 1.1.1，这意味着 $\operatorname{Bun}_G$ 是良的，且
$T^{*}\operatorname{Bun}_G$ 是纯维数 $2\dim\operatorname{Bun}_G$
的局部完全交。

对开子集 $V\subset T^{*}\operatorname{Bun}_G$，以下两条性质等价：
1) $p$ 在 $V$ 上的限制平坦；2) 该限制的纤维维数为
$\dim\operatorname{Bun}_G$。设 $V_{\max}$ 是具有这两条性质的最大 $V$。
$V_{\max}$ 是 $\mathbb{G}_m$-不变的，且按 (38) 有
$V_{\max}\supset p^{-1}(0)$。于是 $V_{\max}=T^{*}\operatorname{Bun}_G$，
这就证明了 2.2.4(iii) 的第一个断言。它蕴含 $p^{\gamma}$ 的像是开的。
另一方面，该像 $\mathbb{G}_m$-不变且包含 $0$。因此 $p^{\gamma}$ 是满射。
证毕。

\SourcePage{49}

由于 $\operatorname{Nilp}$ 包含 $T^{*}\operatorname{Bun}_G$ 的零截面，
(38) 可由不等式
$\dim\operatorname{Nilp}\le\dim\operatorname{Bun}_G$ 推出；后者在
[La88]、[Fal93]、[Gi97] 中作为下述定理的推论得到。

\bdp{2.10.2} \emph{定理.}（[La88]，[Fal93]，[Gi97]）
$\operatorname{Nilp}$ 是迷向的。

\emph{注.}
\begin{enumerate}[label=(\roman*),leftmargin=2.6em]
\item 说明一下：光滑辛簇 $M$ 的子概形 $N$ 称为迷向的，是指 $N$ 的任一
  光滑子簇都是迷向的。可以证明：$N$ 迷向当且仅当 $N_{\mathrm{red}}$
  的非奇异点集迷向。称 $N$ 为 Lagrange 的，是指 $N$ 迷向，且对任意
  $x\in N$ 有 $\dim_x N=1/2\dim_x M$。若 $\mathcal{Y}$ 是光滑代数叠，
  则称子叠 $\mathcal{N}\subset T^{*}\mathcal{Y}$ 为迷向的（相应地，
  Lagrange 的），是指对某个呈现\footnote{$\mathcal{Y}$ 的一个呈现是指
  光滑满态射 $S\to\mathcal{Y}$，其中 $S$ 是概形。}
  $S\to\mathcal{Y}$，有
  $\mathcal{N}\times_{\mathcal{Y}}S\subset
  (T^{*}\mathcal{Y})\times_{\mathcal{Y}}S\subset T^{*}S$
  迷向（相应地，Lagrange）；此时对一切呈现 $S\to\mathcal{Y}$ 也都成立。
\item 定理 2.10.2 在 [Fal93] 与 [Gi97] 中给出的证明都不使用 $g>1$
  的假设，这里 $g$ 是 $X$ 的亏格。若 $g>1$，Faltings 与 Ginzburg
  证明 $\operatorname{Nilp}$ 是 Lagrange 的。其论证已在 2.10.1 中解释：
  (38) 蕴含 $\operatorname{Nilp}$ 的维数纯为
  $\dim\operatorname{Bun}_G$。在 2.10.1 中我们用到了等式
  $\dim\operatorname{Hitch}(X)=\dim\operatorname{Bun}_G$，它只在 $g>1$
  时成立。事实上，即使 $g=0,1$，$\operatorname{Nilp}$ 也是 Lagrange 的
  （见 2.10.4）。
\item 由于 $\operatorname{Nilp}\subset T^{*}\operatorname{Bun}_G$ 是
  Lagrange 的且 $\mathbb{G}_m$-不变，它是一些余法丛的并，这些余法丛对应
  $\operatorname{Bun}_G$ 的某些既约不可约闭子叠。对 $G=PSL_n$，
  Laumon 给出了其中某些子叠的描述（见 [La88] 的 §§3.8--3.9）。
\item Ginzburg 对定理 2.10.2 的证明基于如下解释：用
  $\pi:\operatorname{Bun}_B\to\operatorname{Bun}_G$ 刻画
  $\operatorname{Nilp}$，其中 $B$ 是 $G$ 的 Borel 子群；即若
  $\mathcal{F}\in\operatorname{Bun}_G$，
  $\eta\in T^{*}_{\mathcal{F}}\operatorname{Bun}_G$，则
  $(\mathcal{F},\eta)\in\operatorname{Nilp}$
\SourcePage{50}
  当且仅当存在 $\mathcal{E}\in\pi^{-1}(\mathcal{F})$，使得 $\eta$ 在
  $T^{*}_{\mathcal{E}}\operatorname{Bun}_B$ 中的像等于 $0$。
  借助这一解释，Ginzburg 用辛几何中一个简单而一般的论证证明了
  定理 2.10.2（见 [Gi97] 的 §§6.5）。Faltings 对定理 2.10.2 的证明
  也很漂亮且简短（见 [Fal93] 定理 II.5 证明的前两段）。
\item {}[La88] 中关于 $G=PSL_n$ 的定理 2.10.2 证明在一般情形不成立，
  因为它用到 $\mathfrak{g}=sl_n$ 的如下性质：对每个幂零元
  $A\in\mathfrak{g}$，存在抛物子群 $P\subset G$，使得 $A$ 属于幂幺根
  $U\subset P$ 的李代数，$A$ 的 $P$-轨道在 $\operatorname{Lie}U$ 中
  开，且 $A$ 在 $G$ 中的中心化子含于 $P$。该性质对 $\mathfrak{g}=sl_n$
  成立（例如可取 $P$ 为旗 $0\subset\operatorname{Ker}A
  \subset\operatorname{Ker}A^2\subset\cdots$ 的稳定化子），但对任意
  半单 $\mathfrak{g}$ 不成立（例如当 $\mathfrak{g}=sp_4$、$A\in sp_4$
  是秩 $1$ 的幂零算子时不成立）。
\end{enumerate}

\bdp{2.10.3} 本小节我们「描述」$\operatorname{Nilp}$ 的不可约分支集。

回忆：$\operatorname{Nilp}$ 是二元组 $(\mathcal{F},\eta)$ 构成的叠，
其中 $\mathcal{F}$ 是 $X$ 上 $G$-丛，
$\eta\in H^0(X,\mathfrak{g}_{\mathcal{F}}\otimes\omega_X)
=H^0(X,\mathfrak{g}^{*}_{\mathcal{F}}\otimes\omega_X)$ 幂零。对幂零共轭类
$C\subset\mathfrak{g}$，有局部闭子叠 $\operatorname{Nilp}_C$，
它参数化满足「对一般 $x\in X$ 有 $\eta(x)\in C$」的二元组
$(\mathcal{F},\eta)$。

固定 $e\in C$，把它扩充为 $sl_2$-三元组 $\{e,f,h\}$。设
$\mathfrak{g}^k$ 是 $\mathfrak{g}$ 的降滤链，使得
$[h,\mathfrak{g}^k]\subset\mathfrak{g}^k$，且 $\operatorname{ad}_h$
在 $\mathfrak{g}^k/\mathfrak{g}^{k+1}$ 上以乘以 $k$ 的方式作用。
$\mathfrak{g}^k$ 依赖 $e$，但不依赖 $h$ 与 $f$。置
$\mathfrak{p}=\mathfrak{p}_e:=\mathfrak{g}^0$。$\mathfrak{p}$ 是
$\mathfrak{g}$ 的抛物子代数。设 $P\subset G$ 是相应的子群。有映射
$C\to G/P$，它把 $a\in C$ 对应到抛物子代数 $\mathfrak{p}_a$。其纤维
$\{a\in C\mid\mathfrak{p}_a=\mathfrak{p}\}$（即 $e\in C$ 的 $P$-轨道）
等于 $\mathfrak{g}^2\cap C$；这是 $\mathfrak{g}^2$ 的开子集。称
$\mathfrak{g}^2$ 的元素为一般的，是指它属于 $\mathfrak{g}^2\cap C$。

\SourcePage{51}

设 $(\mathcal{F},\eta)\in\operatorname{Nilp}_C$，
$U:=\{x\in X\mid\eta(x)\in C\}$。$\eta\in
\Gamma(U,C_{\mathcal{F}}\otimes\omega_X)$ 在
$\Gamma(U,(G/P)_{\mathcal{F}})$ 中的像可延拓为 $(G/P)_{\mathcal{F}}$
在 $X$ 上的截面。于是得到 $\mathcal{F}$ 上的 $P$-结构。就这个
$P$-结构而言，$\eta\in H^0(X,\mathfrak{g}^2_{\mathcal{F}}\otimes\omega_X)$，
且对 $x\in U$，$\eta(x)$ 是一般的。

记 $Y_C$ 为二元组 $(\mathcal{F},\eta)$ 构成的叠，其中 $\mathcal{F}$
是 $X$ 上 $P$-丛，$\eta\in
H^0(X,\mathfrak{g}^2_{\mathcal{F}}\otimes\omega_X)$，且对几乎所有
$x\in X$，$\eta(x)$ 是一般的。对 $P$-丛 $\mathcal{F}$，记
$\deg\mathcal{F}\in\operatorname{Hom}(P,\mathbb{G}_m)^*$ 为如下泛函：
它把 $\varphi:P\to\mathbb{G}_m$ 对应到 $\mathcal{F}$ 经 $\varphi$
的前推的次数。$Y_C$ 是开子叠 $Y_C^u$
（$u\in\operatorname{Hom}(P,\mathbb{G}_m)^*$）的无交并，$Y_C^u$
参数化满足 $\deg\mathcal{F}=u$ 的 $(\mathcal{F},\eta)\in Y_C$。
容易证明：对每个 $u\in\operatorname{Hom}(P,\mathbb{G}_m)^*$，
自然态射 $Y_C^u\to\operatorname{Nilp}_C$ 是局部闭嵌入，且子叠
$Y_C^u\subset\operatorname{Nilp}_C$ 构成 $\operatorname{Nilp}_C$
的一个分层。

\emph{引理.}
\begin{enumerate}[label=\arabic*),leftmargin=2.2em]
\item $Y_C^u$ 是光滑等维叠。
  $\dim Y_C^u\le\dim\operatorname{Bun}_G$。
\item 设 $Y_C^*$ 是 $Y_C$ 的维数为 $\dim\operatorname{Bun}_G$ 的连通分支
  之并。则 $Y_C^*$ 是满足下述条件的二元组 $(\mathcal{F},\eta)\in Y_C$
  构成的叠：$\operatorname{ad}_\eta:
  (\mathfrak{g}^{-1}/\mathfrak{g}^0)_{\mathcal{F}}\to
  (\mathfrak{g}^1/\mathfrak{g}^2)_{\mathcal{F}}\otimes\omega_X$
  是同构。
\end{enumerate}

\emph{注.} (38) 由引理推出。

\emph{证明.} 二元组 $(\mathcal{F},\eta)\in Y_C^u$ 的形变理论由复形
$C^{\bullet}$ 的超上同调控制，其中
$C^0=\mathfrak{p}_{\mathcal{F}}=\mathfrak{g}^0_{\mathcal{F}}$，
$C^1=\mathfrak{g}^2_{\mathcal{F}}\otimes\omega_X$，$i\ne 0,1$ 时
$C^i=0$，微分 $d:C^0\to C^1$ 等于 $\operatorname{ad}_\eta$。由于
$\operatorname{Coker}d$ 有有限支撑，
$\mathbb{H}^2(X,C^{\bullet})=0$。因此 $Y_C$ 光滑，且
\[
\begin{aligned}
\dim_{(\mathcal{F},\eta)}Y_C
&=\chi(\mathfrak{g}^2_{\mathcal{F}}\otimes\omega_X)
  -\chi(\mathfrak{g}^0_{\mathcal{F}})
 =-\chi(\mathfrak{g}_{\mathcal{F}}/\mathfrak{g}^{-1}_{\mathcal{F}})
  -\chi(\mathfrak{g}^0_{\mathcal{F}})\\
&=-\chi(\mathfrak{g}_{\mathcal{F}})
  +\chi(\mathfrak{g}^{-1}_{\mathcal{F}}/\mathfrak{g}^0_{\mathcal{F}})
 =\dim\operatorname{Bun}_G
  +\chi(\mathfrak{g}^{-1}_{\mathcal{F}}/\mathfrak{g}^0_{\mathcal{F}}).
\end{aligned}
\]
显然 $\chi(\mathfrak{g}^{-1}_{\mathcal{F}}/\mathfrak{g}^0_{\mathcal{F}})$
只依赖 $u=\deg\mathcal{F}$。态射 $\operatorname{ad}_\eta:
\mathfrak{g}^{-1}_{\mathcal{F}}/\mathfrak{g}^0_{\mathcal{F}}\to
(\mathfrak{g}^1/\mathfrak{g}^2)_{\mathcal{F}}\otimes\omega_X$ 是单射，
其余核 $\mathcal{A}$ 有有限支撑。于是
$2\chi(\mathfrak{g}^{-1}_{\mathcal{F}}/\mathfrak{g}^0_{\mathcal{F}})
=\chi(\mathfrak{g}^{-1}_{\mathcal{F}}/\mathfrak{g}^0_{\mathcal{F}})
-\chi((\mathfrak{g}^1/\mathfrak{g}^2)_{\mathcal{F}}\otimes\omega_X)
=-\chi(\mathcal{A})\le 0$，且
$\chi(\mathfrak{g}^{-1}_{\mathcal{F}}/\mathfrak{g}^0_{\mathcal{F}})=0$
当且仅当 $\mathcal{A}=0$。\hfill$\Box$

\SourcePage{52}

由于 $\operatorname{Nilp}$ 的维数纯为 $\dim\operatorname{Bun}_G$，
引理蕴含：$\operatorname{Nilp}$ 的不可约分支由
$\bigsqcup_C\pi_0(Y_C^*)$ 参数化。

$\pi_0(Y_C^*)$ 可等同于一个更简单的叠 $M_C$ 的 $\pi_0$，后者定义如下。
置 $L=P/U$，其中 $U$ 是 $P$ 的幂幺根。$L$ 作用在
$V:=\mathfrak{g}^2/\mathfrak{g}^3$ 上。记 $D_i$ 为满足下述条件的
$a\in V$ 的集合：$(\operatorname{ad}_a)^i:
\mathfrak{g}^{-i}/\mathfrak{g}^{-i+1}\to
\mathfrak{g}^i/\mathfrak{g}^{i+1}$ 的行列式为 $0$。
$D_i\subset V$ 是纯余维 $1$ 的 $L$-不变闭子集。$\mathfrak{g}^2$
的元素一般，当且仅当它在 $V$ 中的像不属于 $D_2$。因此对所有 $i$ 有
$D_i\subset D_2$。记 $M_C$ 为二元组 $(\mathcal{F},\eta)$ 构成的叠，
其中 $\mathcal{F}$ 是 $X$ 上 $L$-丛，
$\eta\in H^0(X,V_{\mathcal{F}}\otimes\omega_X)$ 满足：对一切
$x\in X$ 有 $\eta(x)\notin D_1$，且对一般的 $x\in X$ 有
$\eta(x)\notin D_2$。容易看出，自然态射 $Y_C^*\to M_C$ 诱导双射
$\pi_0(Y_C^*)\to\pi_0(M_C)$。

于是 $\operatorname{Nilp}$ 的不可约分支由
$\bigsqcup_C\pi_0(M_C)$ 参数化。但愿 $\pi_0(M_C)$ 能用与 $C$ 和
$X$ 相伴的「标准」对象来描述……

\emph{注.} 若 $G=PSL_n$，则对每个幂零共轭类 $C\subset sl_n$，
$\operatorname{Nilp}_C$ 的维数纯为 $\dim\operatorname{Bun}_G$
（见 [La88]）。这并非总成立。例如取 $G=Sp_4$，$C$ 为 $sp_4$ 中秩
$1$ 的幂零矩阵之集。事实上，设 $(\mathcal{F},\eta)\in Y_C$ 使得
$\eta\in H^0(X,\mathfrak{g}^2_{\mathcal{F}}\otimes\omega_X)$ 只有单零点。
则容易证明：态射 $Y_C\to\operatorname{Nilp}_C$ 在
$(\mathcal{F},\eta)$ 的邻域内是开嵌入。另一方面，由上述引理，若
$\eta$ 有零点，则 $Y_C$ 在 $(\mathcal{F},\eta)$ 处的维数小于
$\operatorname{Bun}_G$。

\bdp{2.10.4} \emph{定理.} $\operatorname{Nilp}$ 是 Lagrange 的。

本定理不假设 $g>1$。

\emph{证明.} 按 2.10.2 的注 (ii)，只需证明当 $g\le 1$ 时
$\operatorname{Nilp}$ 的维数纯为 $\dim\operatorname{Bun}_G$。

1) 设 $g=0$。则 $\operatorname{Nilp}=T^{*}\operatorname{Bun}_G$。
$\operatorname{Bun}_G$ 的拟紧开子叠可表为 $H\backslash M$，其中 $M$
是光滑簇，$H$ 是
% [Tutor Note: 原书 p.52 在此处断页，句子未完，下接 p.53。]
